\chapter{Introduction}

\section{Motivation}

According to the World Health Organization, more than 5\% of the world's
population lives with disabling hearing loss. For many of these people, a
signed language—rather than a written or spoken one—is their first and most
natural means of communication. American Sign Language (ASL) is a complete,
grammatically rich language used by a large community across North America and
beyond. Despite this, communication between signers and the hearing majority
remains difficult in everyday situations such as shops, clinics, and public
offices, because qualified interpreters are scarce and cannot be present on
demand.

Technology offers a path to narrow this gap. If a wearable device could
recognise signed gestures and render them as speech, a signer could be
understood by a hearing person who knows no sign language, with no third party
required. Such a device must be \emph{portable}, \emph{low-latency},
\emph{affordable}, and ideally \emph{self-contained}, so that it works anywhere
without a network connection or a fixed camera installation. These constraints
motivate the design presented in this thesis: a glove that senses hand and
finger motion with small inertial sensors and performs the entire recognition
pipeline on-device.

\section{Problem Statement}

The core problem is to translate a vocabulary of ASL gestures into text and
speech using only wearable inertial sensing, subject to the practical
constraints of a battery- or USB-powered wearable. Concretely, the system must:

\begin{itemize}
    \item sense the configuration and motion of the five fingers with
    sufficient fidelity to distinguish a useful vocabulary of letters and
    words;
    \item transmit that data wirelessly and reliably from the hand to a
    processing device;
    \item classify gestures accurately and with low latency;
    \item operate without a camera, a cloud service, or a wired tether during
    use; and
    \item remain robust to the practical realities of a wearable—imperfect
    re-calibration between uses, limited power, and the heat generated by the
    on-hand electronics.
\end{itemize}

\section{Objectives and Scope}

The objectives of this project are:

\begin{enumerate}
    \item To design and build a five-finger sensing glove based on commodity
    inertial measurement units and microcontrollers.
    \item To establish a reliable wireless data path from the glove to a mobile
    phone.
    \item To collect a labelled dataset of ASL gestures and to train a compact
    neural network suitable for on-device inference.
    \item To develop a mobile application that performs inference and produces
    spoken output.
    \item To evaluate the accuracy, latency, and practical robustness of the
    complete system, and to document the engineering decisions and trade-offs
    encountered.
\end{enumerate}

The scope is limited to a single hand, a fixed vocabulary of fingerspelling
letters and a set of common one-handed word signs, and recognition of
deliberate, segmented gestures (one gesture captured per request) rather than
continuous, unsegmented signing.

\section{Contributions}

The principal contributions of this work are:

\begin{itemize}
    \item A \textbf{no-hub, master/slave wearable architecture} in which one
    finger module aggregates the data of the other four and bridges directly
    to a phone, eliminating the separate wrist-mounted receiver and the
    external antennas of earlier iterations.
    \item An empirical study of \textbf{feature representation} for
    inertial sign recognition, showing why orientation-normalised
    (``relative'') features and absolute features each succeed or fail
    depending on whether a gesture is static or dynamic, and arriving at a
    \emph{dynamic-gesture, absolute-feature} formulation that is robust to
    imperfect re-calibration.
    \item A \textbf{train/inference-parity pipeline} in which the exact
    windowing, feature transform, and normalisation used during data
    collection are reproduced on the phone, together with a quantised model
    that runs in roughly one millisecond on-device.
    \item A documented account of the \textbf{system-level engineering
    trade-offs}—radio coexistence, USB throughput, power integrity, and
    thermal behaviour—that are rarely reported but decisive in practice.
\end{itemize}

\section{Thesis Outline}

Chapter~2 reviews sign-language recognition approaches, inertial sensing, and
edge machine learning. Chapter~3 describes the system requirements and the
hardware architecture, including the evolution of the design. Chapter~4 details
the firmware and the sensor signal-processing chain. Chapter~5 presents the
data-acquisition pipeline, the feature-representation study, and the neural
network. Chapter~6 describes the mobile application and the on-device inference
path. Chapter~7 reports and discusses the evaluation results. Chapter~8
concludes and outlines future work.
